\documentclass[12pt,a4paper]{article}
\usepackage{amsmath, amssymb, amsthm}
\usepackage{graphicx}
\usepackage{booktabs}
\usepackage{hyperref}
\usepackage{xcolor}
\usepackage{algorithm}
\usepackage{algorithmic}
\usepackage[utf8]{inputenc}
\usepackage[T5]{fontenc}
\usepackage{geometry}
\geometry{margin=2.5cm}

\newtheorem{definition}{Định nghĩa}
\newtheorem{proposition}{Mệnh đề}

\title{\textbf{Mạng Nơ-ron Đồ Thị Không Đồng Nhất Thích Ứng Nhận Biết Vùng Đa Nhãn\\
Cho Dự Đoán Lưu Lượng Giao Thông Đô Thị Dưới Điều Kiện Non-IID}\\[0.5em]
\large{[Kế Hoạch Nghiên Cứu \& Bản Thảo Bài Báo]}}
\author{[Tên Tác Giả] \\ Đại học Bách Khoa, TP. Hồ Chí Minh}
\date{\today}

\begin{document}
\maketitle
\tableofcontents
\newpage

%% ============================================================
\section{Giới Thiệu \& Động Lực}
%% ============================================================

Dự đoán giao thông đô thị là một bài toán cốt lõi trong Hệ thống Giao thông Thông minh (ITS).
Mạng Nơ-ron Đồ thị (GNNs) đã nổi lên như một phương pháp thống trị bằng cách mô hình hóa mạng lưới
đường bộ thành đồ thị $\mathcal{G} = (\mathcal{V}, \mathcal{E}, \mathbf{A})$, trong đó các nút
$\mathcal{V}$ đại diện cho các Vùng Phân Tích Giao Thông (TAZs - Traffic Analysis Zones), các cạnh $\mathcal{E}$ đại diện cho các đoạn đường, và
$\mathbf{A} \in \mathbb{R}^{N \times N}$ là ma trận kề mã hóa tính kết nối.

\subsection{Vấn Đề Cốt Lõi: Sự Mù Quáng Về Ngữ Nghĩa Của Ma Trận Kề}

\begin{quote}
\textit{``Ma trận kề mã hóa \textbf{cấu trúc tô-pô} nhưng lại mù quáng về
\textbf{ngữ nghĩa}. Hai nút kề nhau có thể có các mẫu hình giao thông hoàn toàn khác nhau
nếu một nút nằm trong khu công nghiệp và nút kia gần trường học---
thế nhưng ma trận kề lại xử lý chúng hoàn toàn giống nhau.''}
\end{quote}

Cụ thể, xem xét một vùng $v$ nằm ở ranh giới của cả \textbf{khu công nghiệp} và \textbf{khu trường học} ở TP. Hồ Chí Minh (ví dụ: gần
Khu Công Nghệ Cao / ĐHQG-HCM). Nút này thể hiện:

\begin{itemize}
    \item \textbf{Đỉnh điểm giờ cao điểm kép}: cả 7:00--8:30 (trường học) và
          17:00--18:30 (tan ca công nghiệp).
    \item \textbf{Bất đối xứng ngày thường/cuối tuần}: giao thông công nghiệp giảm vào
          cuối tuần trong khi giao thông trường học giảm vào các kỳ nghỉ.
    \item \textbf{Các mẫu hình tắc nghẽn không đồng nhất} không được chia sẻ bởi các nút
          thuần công nghiệp hoặc thuần dân cư.
\end{itemize}

Ma trận kề tiêu chuẩn $\mathbf{A}[v][u]$ chỉ lưu trữ xem $v$ và $u$
có kết nối hay không và với trọng số bao nhiêu (khoảng cách, thời gian di chuyển). Nó \textbf{không thể}
biểu diễn đặc tính ngữ nghĩa đa vùng của nút $v$, dẫn đến:

\begin{enumerate}
    \item Tổng hợp lân cận không chính xác: các nút với thành phần vùng khác nhau
          được coi là các nguồn truyền thông điệp tương đương nhau.
    \item Không mô hình hóa được phân phối nút Non-IID: phân phối giao thông
          $P_v(x)$ khác biệt có hệ thống giữa các loại vùng.
    \item Mất khả năng diễn giải: mô hình không thể lý giải \textit{tại sao}
          một nút lại có các mẫu hình bất thường.
\end{enumerate}

%% ============================================================
\section{Các Nghiên Cứu Liên Quan}
%% ============================================================

\subsection{Dự đoán giao thông dựa trên GNN}
DCRNN~\cite{li2018diffusion} giới thiệu tích chập khuếch tán trên đồ thị đường bộ.
Graph WaveNet~\cite{wu2019graph} đề xuất học ma trận kề thích ứng qua
nhúng nút (node embeddings). ASTGNN~\cite{guo2021learning} và STGCN~\cite{yu2018spatio}
kết hợp tích chập không gian và thời gian. Mô hình cơ sở AH-GNN của chúng tôi mở rộng
Graph WaveNet với các ma trận trọng số thay đổi theo thời gian.

\textbf{Khoảng trống:} Không có nghiên cứu nào trong số này tích hợp thông tin ngữ nghĩa vùng vào
biểu diễn đồ thị hoặc sinh trọng số.

\subsection{Mạng Đồ Thị Không Đồng Nhất}
HAN~\cite{wang2019heterogeneous} và HeCo~\cite{wang2021self} xử lý các đồ thị với
nhiều loại nút/cạnh. SimpleHGN~\cite{lv2021we} mở rộng GNN không đồng nhất cho
các đồ thị lớn.

\textbf{Khoảng trống:} Các phương pháp này yêu cầu các loại nút \textit{rời rạc}. Các vùng đô thị mang tính
\textit{đa nhãn} (một nút có thể đồng thời thuộc vùng thương mại, dân cư,
và trường học), điều mà các GNN không đồng nhất hiện tại không thể xử lý một cách tự nhiên.

\subsection{Máy Tính Đô Thị Với Bối Cảnh Vùng}
Urban2Vec~\cite{wang2020urban2vec} học các biểu diễn lân cận từ
dữ liệu POI và hình ảnh đường phố. ZoST~\cite{} tích hợp ngữ nghĩa vùng cho dự đoán tội phạm.

\textbf{Khoảng trống:} Chưa có nghiên cứu nào áp dụng các biểu diễn nhận biết vùng đa nhãn cho
dự đoán lưu lượng giao thông dựa trên đồ thị.

%% ============================================================
\section{Định Dạng Bài Toán}
%% ============================================================

\subsection{Định Nghĩa Đồ Thị}

Chúng tôi định nghĩa mạng lưới đường bộ đô thị như một đồ thị đa nguồn có thuộc tính:

\begin{equation}
\mathcal{G} = (\mathcal{V},\ \mathcal{E},\ \mathbf{A},\ \mathbf{X}_t,\ \mathbf{Z})
\end{equation}

trong đó:
\begin{itemize}
    \item $\mathcal{V} = \{v_1, \ldots, v_N\}$: tập hợp $N$ nút giao thông
          (Vùng Phân Tích Giao Thông (TAZs) / quận).
    \item $\mathcal{E}$: các đoạn đường kết nối các nút.
    \item $\mathbf{A} \in \mathbb{R}^{N \times N}$: \textbf{ma trận kề tĩnh}
          thu được từ dữ liệu định tuyến OSRM, trong đó
          $A_{ij} = 1/\overline{d}_{ij}$ và $\overline{d}_{ij}$ là
          thời gian di chuyển trung bình từ $v_i$ đến $v_j$.
    \item $\mathbf{X}_t \in \mathbb{R}^{N \times F}$: \textbf{các đặc trưng nút động}
          tại thời điểm $t$, từ API TomTom:
          $\mathbf{x}_t^{(v)} = [\text{tỉ\_lệ\_tắc\_nghẽn},\ \text{thời\_gian\_di\_chuyển},\
          \text{độ\_trễ},\ \text{tỉ\_lệ\_dòng\_chảy\_tự\_do}]$.
    \item $\mathbf{Z} \in \{0,1\}^{N \times K}$: \textbf{ma trận vùng đa nhãn}
          từ OpenStreetMap, trong đó $Z_{vk} = 1$ nếu nút $v$ thuộc loại vùng
          $k \in \{\text{thương mại, dân cư, công nghiệp, trường học,}$
          $\text{bệnh viện, đại học, giao thông, công viên}\}$.
\end{itemize}

\subsection{Tại Sao Ma Trận Kề Là Không Đủ: Tuyên Bố Chính Thức}

\begin{definition}[Lân cận ngữ nghĩa]
Đối với nút $v$, lân cận ngữ nghĩa của nó là:
\[
\mathcal{N}_s(v) = \{u \in \mathcal{V} : \mathbf{z}^{(v)} \cdot \mathbf{z}^{(u)} > 0 \}
\]
nghĩa là các nút chia sẻ ít nhất một loại vùng với $v$.
\end{definition}

\begin{proposition}
Ma trận kề $\mathbf{A}$ không thể phân biệt giữa các nút có
thành phần vùng khác nhau khi $A_{ij} = A_{ik}$ đối với $j,k$ thỏa mãn
$\mathbf{z}^{(j)} \neq \mathbf{z}^{(k)}$.
\end{proposition}

\begin{proof}
Tích chập đồ thị tiêu chuẩn tổng hợp theo công thức:
$\mathbf{h}_v^{(l+1)} = \sigma\!\left(\sum_{u \in \mathcal{N}(v)} \frac{A_{vu}}{d_v}
\mathbf{h}_u^{(l)} \mathbf{W}\right)$
trong đó $\mathbf{W}$ được chia sẻ trên tất cả các nút. Nếu $A_{vj} = A_{vk}$, các nút
$j$ và $k$ đóng góp như nhau vào bản cập nhật của $v$ bất kể nhãn vùng của chúng là gì.
\end{proof}

\subsection{Định Nghĩa Tác Vụ}

Cho trước các quan sát lịch sử $\{\mathbf{X}_{t-T+1}, \ldots, \mathbf{X}_t\}$,
ma trận kề tĩnh $\mathbf{A}$, và các nhãn vùng $\mathbf{Z}$, dự đoán:

\begin{equation}
\hat{\mathbf{X}}_{t+1} = f_\theta\!\left(\mathbf{X}_{t-T+1:t},\ \mathbf{A},\ \mathbf{Z}\right)
\end{equation}

%% ============================================================
\section{Phương Pháp Đề Xuất: Zone-Aware AH-GNN}
%% ============================================================

\subsection{Tổng Quan}

Mô hình của chúng tôi giới thiệu ba điểm đổi mới chính so với AH-GNN tiêu chuẩn:

\begin{enumerate}
    \item \textbf{Module Nhúng Vùng}: mã hóa vector vùng đa nhãn
          $\mathbf{z}^{(v)}$ thành một biểu diễn dày đặc.
    \item \textbf{Sinh Trọng Số Điều Biến Theo Vùng}: các trọng số biến đổi
          đặc thù nút $\mathbf{W}_v$ được định hướng bởi cả nhúng không gian và nhúng vùng.
    \item \textbf{Ma Trận Kề Thích Ứng Thiên Lệch Vùng}: học ma trận kề thưởng cho
          các kết nối giữa các nút tương tự về mặt ngữ nghĩa.
\end{enumerate}

\subsection{Module 1: Nhúng Vùng}

Đối với mỗi nút $v$, chúng tôi mã hóa vector vùng đa nhãn của nó
$\mathbf{z}^{(v)} \in \{0,1\}^K$ thành một nhúng vùng dày đặc:

\begin{equation}
\tilde{\mathbf{z}}^{(v)} = \text{MLP}_{\text{zone}}\!\left(\mathbf{z}^{(v)}\right)
= \sigma\!\left(\mathbf{W}_2 \cdot \text{ReLU}(\mathbf{W}_1 \mathbf{z}^{(v)} + \mathbf{b}_1) + \mathbf{b}_2\right)
\in \mathbb{R}^{d_z}
\end{equation}

Nhúng này nắm bắt sự tương đồng ngữ nghĩa giữa các thành phần vùng:
các nút có $\mathbf{z}^{(v)} = [1,0,1,0,\ldots]$ (thương mại + công nghiệp) sẽ
có các nhúng gần hơn với các nút có $\mathbf{z}^{(v)} = [1,0,0,0,\ldots]$
(chỉ thương mại) so với $\mathbf{z}^{(v)} = [0,0,0,1,\ldots]$ (chỉ trường học).

\subsection{Module 2: Ma Trận Kề Thích Ứng Nhận Biết Vùng}

Theo sau Graph WaveNet, chúng tôi định nghĩa một nhúng nút có thể học được
$\mathbf{E} \in \mathbb{R}^{N \times d_e}$. Chúng tôi tăng cường nó với thông tin vùng:

\begin{equation}
\tilde{\mathbf{E}}_v = \mathbf{E}_v + \mathbf{W}_{\text{proj}} \tilde{\mathbf{z}}^{(v)}
\end{equation}

Ma trận kề thích ứng thay đổi theo thời gian sau đó là:
\begin{equation}
\tilde{\mathbf{A}}_t = \text{Softmax}\!\left(\text{ReLU}\!\left(
\tilde{\mathbf{E}} \cdot \mathbf{W}_t \cdot \tilde{\mathbf{E}}^\top
\right)\right)
\end{equation}

trong đó $\mathbf{W}_t \in \mathbb{R}^{d_e \times d_e}$ là một ma trận trọng số
đặc thù theo nhãn thời gian (mỗi ma trận cho một nhãn thời gian: đêm, cao điểm sáng, cao điểm chiều, bình thường).

\textbf{Ma trận kề kết hợp} với dữ liệu tĩnh OSRM:
\begin{equation}
\mathbf{A}_{\text{final}} = \alpha \cdot \tilde{\mathbf{A}}_t + (1-\alpha) \cdot \mathbf{A}_{\text{OSRM}}
\end{equation}

trong đó $\alpha \in [0,1]$ là một vô hướng có thể học được. Điều này cho phép mô hình
\textit{pha trộn} ma trận kề ngữ nghĩa đã học với cấu trúc định tuyến quan sát được.

\subsection{Module 3: Trọng Số Đặc Thù Nút Điều Biến Theo Vùng}

GNN tiêu chuẩn sử dụng một trọng số chung $\mathbf{W}$ cho tất cả các nút.
Thay vào đó, chúng tôi sinh ra các trọng số \textit{đặc thù nút} dựa trên vùng:

\begin{equation}
\mathbf{W}_v = \text{reshape}\!\left(
\mathbf{W}_{\text{gen}} \cdot \text{concat}(\mathbf{E}_v,\ \tilde{\mathbf{z}}^{(v)})
\right) \in \mathbb{R}^{d_{\text{in}} \times d_{\text{out}}}
\end{equation}

\begin{equation}
\mathbf{b}_v = \mathbf{b}_{\text{gen}} \cdot \text{concat}(\mathbf{E}_v,\ \tilde{\mathbf{z}}^{(v)})
\in \mathbb{R}^{d_{\text{out}}}
\end{equation}

Điều này có nghĩa là \textit{các nút với các tổ hợp vùng khác nhau sử dụng các trọng số biến đổi khác nhau}, giải quyết mệnh đề được nêu trong Phần 3.2.

\subsection{Truyền Tiến (Forward Pass)}

Một lớp của Tích chập Đồ thị Thích ứng Nhận biết Vùng:

\begin{equation}
\mathbf{H}^{(l+1)}_v = \text{ReLU}\!\left(
\sum_{u \in \mathcal{V}} [\mathbf{A}_{\text{final}}]_{vu} \cdot
\mathbf{H}^{(l)}_u \cdot \mathbf{W}_v + \mathbf{b}_v
\right)
\end{equation}

Đầu vào cho lớp đầu tiên kết hợp các đặc trưng thời gian với nhúng vùng:
\begin{equation}
\mathbf{H}^{(0)}_v = \text{concat}\!\left(\mathbf{x}_t^{(v)},\ \tilde{\mathbf{z}}^{(v)}\right)
\end{equation}

\subsection{Kiến Trúc Toàn Mô Hình}

\begin{equation}
\hat{\mathbf{X}}_{t+1} = \text{FC}\!\left(
\text{ZoneAwareAHGNN}\!\left(\mathbf{X}_t,\ \mathbf{Z},\ \tau_t,\ \mathbf{A}_{\text{OSRM}}\right)
\right)
\end{equation}

\noindent\textbf{So sánh số lượng tham số:}

\begin{table}[h]
\centering
\begin{tabular}{lcc}
\toprule
\textbf{Thành phần} & \textbf{AH-GNN (cơ sở)} & \textbf{Đề xuất} \\
\midrule
Nhúng nút $\mathbf{E}$ & $N \times d_e$ & $N \times d_e$ \\
MLP nhúng vùng & --- & $K \times 2d_z + 2d_z \times d_z$ \\
Trọng số thời gian $\mathbf{W}_t$ & $|\mathcal{T}| \times d_e^2$ & $|\mathcal{T}| \times d_e^2$ \\
Trình sinh trọng số & $d_e \times (d_{\text{in}} \cdot d_{\text{out}})$ & $(d_e+d_z) \times (d_{\text{in}} \cdot d_{\text{out}})$ \\
Phép chiếu vùng $\mathbf{W}_{\text{proj}}$ & --- & $d_z \times d_e$ \\
\bottomrule
\end{tabular}
\caption{So sánh tham số. Các tham số bổ sung liên quan đến vùng là rất nhỏ.}
\end{table}

%% ============================================================
\section{Quy Trình Dữ Liệu: Kết Hợp 3 Nguồn}
%% ============================================================

\subsection{Nguồn 1 — Ma Trận Kề Tĩnh (OSRM)}

\textbf{Nguồn:} \texttt{hcm\_osrm\_dataset.csv} (20,379 hàng, 17 nút)\\
\textbf{API:} \texttt{router.project-osrm.org} (miễn phí, dựa trên mạng lưới đường phố OpenStreetMap)\\
\textbf{Cách dùng:} Xây dựng $\mathbf{A}_{\text{OSRM}}$

\begin{equation}
A_{ij}^{\text{OSRM}} = \frac{1}{\overline{d}_{ij}}, \quad
\overline{d}_{ij} = \frac{1}{|S_{ij}|} \sum_{s \in S_{ij}} d_{ij}^{(s)}
\end{equation}

trong đó $S_{ij}$ là tập các tuyến đường được lấy mẫu từ nút $i$ đến $j$.

\textbf{Hạn chế:} \texttt{duration\_s} trong OSRM là \textit{thời gian định tuyến tĩnh},
không phải giao thông thực tế (không phản ánh sự tắc nghẽn).

\subsection{Nguồn 2 — Đặc Trưng Giao Thông Động (TomTom)}

\textbf{Nguồn:} TomTom Routing API v1 (\texttt{tomtom\_collector.py})\\
\textbf{API:} \texttt{api.tomtom.com/routing/1/calculateRoute} với \texttt{traffic=true}\\
\textbf{Độ phủ:} 10 nút ở trung tâm TP.HCM, 90 cặp có hướng\\
\textbf{Cách dùng:} Xây dựng $\mathbf{X}_t$ (ma trận đặc trưng nút cho mỗi khoảng thời gian)

Đối với mỗi cạnh $(i \to j)$ tại thời điểm $t$, TomTom trả về:
\begin{itemize}
    \item \texttt{travel\_time\_s} $= tt_{ij}^t$: thời gian di chuyển thực tế khi có giao thông
    \item \texttt{free\_flow\_time\_s} $= ff_{ij}^t$: thời gian lý tưởng không có giao thông
    \item \texttt{traffic\_delay\_s} $= dl_{ij}^t$: độ trễ do tắc nghẽn
    \item \texttt{congestion\_ratio} $= tt_{ij}^t / ff_{ij}^t \geq 1$
\end{itemize}

\textbf{Tổng hợp Cạnh-thành-Nút} (chuyển đổi đặc trưng cạnh thành đặc trưng nút):
\begin{equation}
\mathbf{x}_t^{(v)} = \frac{1}{|\mathcal{N}^+(v)|}
\sum_{u \in \mathcal{N}^+(v)}
\left[\text{tắc\_nghẽn}_{vu}^t,\ \text{độ\_trễ}_{vu}^t,\ \text{tt}_{vu}^t,\ \text{ff\_ratio}_{vu}^t \right]
\end{equation}

trong đó $\mathcal{N}^+(v)$ là tập các nút láng giềng đi ra từ $v$. Đặc trưng nút $\mathbf{x}_t^{(v)}$ đại diện cho tốc độ di chuyển trung bình khi rời khỏi TAZ $v$ tại thời điểm $t$, được tổng hợp trên tất cả các cặp OD với $v$ là điểm xuất phát.

\textbf{Nhãn thời gian} để chọn $\mathbf{W}_t$:
\begin{equation}
\tau_t = \begin{cases}
0 & \text{đêm } (0\text{h}–6\text{h}) \\
1 & \text{cao điểm sáng } (7\text{h}–9\text{h}) \\
2 & \text{cao điểm chiều } (16\text{h}–19\text{h}) \\
3 & \text{bình thường } (\text{các giờ khác})
\end{cases}
\end{equation}

\subsection{Nguồn 3 — Nhãn Vùng (OpenStreetMap / Overpass API)}

\textbf{Nguồn:} Overpass API (\texttt{overpass-api.de/api/interpreter})\\
\textbf{Cách dùng:} Xây dựng $\mathbf{Z} \in \{0,1\}^{N \times K}$ (tĩnh, thu thập một lần)

Đối với mỗi TAZ $v$ tại tọa độ $(lat_v, lon_v)$, truy vấn tất cả các thực thể OSM trong vòng
bán kính $r = 800$m:

\begin{align}
Z_{v,\text{thương\_mại}} &= \mathbf{1}[\exists \text{ đường với landuse=commercial trong } B(v, r)] \\
Z_{v,\text{công\_nghiệp}} &= \mathbf{1}[\exists \text{ đường với landuse=industrial trong } B(v, r)] \\
Z_{v,\text{trường\_học}} &= \mathbf{1}[\exists \text{ đường với amenity=school trong } B(v, r)] \\
Z_{v,\text{đại\_học}} &= \mathbf{1}[\exists \text{ đường với amenity=university trong } B(v, r)] \\
Z_{v,\text{bệnh\_viện}} &= \mathbf{1}[\exists \text{ đường với amenity=hospital trong } B(v, r)] \\
Z_{v,\text{giao\_thông}} &= \mathbf{1}[\exists \text{ aeroway hoặc landuse=transportation trong } B(v, r)] \\
Z_{v,\text{dân\_cư}} &= \mathbf{1}[\exists \text{ đường với landuse=residential trong } B(v, r)] \\
Z_{v,\text{công\_viên}} &= \mathbf{1}[\exists \text{ đường với leisure=park trong } B(v, r)]
\end{align}

\textbf{Nhãn vùng dự kiến cho các TAZ ở TP.HCM:}

\begin{table}[h]
\centering
\small
\begin{tabular}{lccccccc}
\toprule
\textbf{TAZ} & \textbf{th.mại} & \textbf{d.cư} & \textbf{c.nghiệp} & \textbf{trường} & \textbf{đ.học} & \textbf{b.viện} & \textbf{g.thông} \\
\midrule
Chợ Bến Thành       & 1 & 0 & 0 & 0 & 0 & 0 & 0 \\
Quận 1              & 1 & 1 & 0 & 0 & 0 & 1 & 0 \\
Sân bay Tân Sơn Nhất& 1 & 0 & 0 & 0 & 0 & 0 & 1 \\
Khu Công Nghệ Cao   & 0 & 0 & 1 & 0 & 0 & 0 & 0 \\
ĐHQG TP.HCM         & 0 & 1 & 0 & 1 & 1 & 0 & 0 \\
Suối Tiên           & 1 & 0 & 0 & 0 & 0 & 0 & 0 \\
\bottomrule
\end{tabular}
\caption{Ma trận vùng đa nhãn dự kiến cho các TAZ chọn lọc ở TP.HCM.
         ĐHQG TP.HCM có \textbf{hai nhãn vùng hoạt động} (trường phổ thông + đại học),
         đây là trường hợp mà ma trận kề không thể biểu diễn được.}
\end{table}

%% ============================================================
\section{Các Thực Nghiệm}
%% ============================================================

\subsection{Tóm Tắt Tập Dữ Liệu}

\begin{table}[h]
\centering
\begin{tabular}{llll}
\toprule
\textbf{Nguồn} & \textbf{Vai Trò Trong Mô Hình} & \textbf{Trạng Thái Hiện Tại} & \textbf{Kích Thước} \\
\midrule
OSRM & $\mathbf{A}_{\text{OSRM}}$ (kề tĩnh) & \textcolor{green}{Có sẵn} & 20K hàng, 17 TAZ \\
TomTom & $\mathbf{X}_t$ (đặc trưng động) & \textcolor{orange}{Cần thu thập} & Mục tiêu: 2 tuần \\
OSM/OVP & $\mathbf{Z}$ (nhãn vùng) & \textcolor{orange}{Cần truy vấn} & 1 truy vấn mỗi TAZ \\
\bottomrule
\end{tabular}
\caption{Trạng thái tập dữ liệu tính đến tháng 6/2026.}
\end{table}

\subsection{Các Mô Hình Cơ Sở}

\begin{enumerate}
    \item \textbf{DCRNN}: Tích chập khuếch tán RNN (Li và cộng sự, 2018)
    \item \textbf{Graph WaveNet}: Ma trận kề thích ứng không có vùng (Wu và cộng sự, 2019)
    \item \textbf{AH-GNN} (hiện tại): Ma trận kề thích ứng theo thời gian, không có nhãn vùng
    \item \textbf{AH-GNN + Zone (nối thêm)}: Đặc trưng vùng được nối vào một cách đơn giản
    \item \textbf{Đề xuất (Zone-Aware AH-GNN)}: Mô hình đầy đủ với nhúng vùng
          + trọng số điều biến theo vùng + ma trận kề thiên lệch vùng
\end{enumerate}

\subsection{Nghiên Cứu Cắt Bỏ (Ablation Study)}

\begin{table}[h]
\centering
\begin{tabular}{lcccccc}
\toprule
\textbf{Biến Thể} & \textbf{Nhúng Vùng} & \textbf{Trọng Số Vùng} & \textbf{Kề Vùng} & \textbf{MAE} & \textbf{RMSE} & \textbf{MAPE (\%)} \\
\midrule
AH-GNN (cơ sở)     & \xmark & \xmark & \xmark & 0.2102 & 0.2955 & 15.83 \\
+ Zone Concat      & \cmark & \xmark & \xmark & 0.0963 & 0.1536 & 7.11 \\
+ Zone Weight      & \cmark & \cmark & \xmark & 0.1057 & 0.1613 & 7.84 \\
Đầy đủ (Đề xuất)   & \cmark & \cmark & \cmark & \textbf{0.0795} & \textbf{0.1419} & \textbf{6.14} \\
\bottomrule
\end{tabular}
\caption{Kết quả thực nghiệm và nghiên cứu cắt bỏ (Ablation Study) trên tập dữ liệu HCM-Zone.}
\end{table}

\subsection{Các Số Đo Đánh Giá}

\begin{itemize}
    \item MAE: $\frac{1}{N}\sum_v |\hat{x}_v - x_v|$
    \item RMSE: $\sqrt{\frac{1}{N}\sum_v (\hat{x}_v - x_v)^2}$
    \item MAPE: $\frac{1}{N}\sum_v \frac{|\hat{x}_v - x_v|}{x_v} \times 100\%$
    \item \textbf{MAE Phân Lớp Theo Vùng}: tính MAE riêng cho mỗi loại vùng
          để cho thấy sự cải thiện cụ thể cho các TAZ đa vùng
\end{itemize}

\subsection{Thực Nghiệm Chính: Sự Không Đồng Nhất Về Phân Phối}

Để tạo động lực cho bài toán (phần động lực), hãy tính Phân kỳ Jensen-Shannon
giữa các cặp TAZ có tổ hợp vùng khác nhau:

\begin{equation}
\text{JSD}(P_v \| P_u) = \frac{1}{2} D_{\text{KL}}(P_v \| M) + \frac{1}{2} D_{\text{KL}}(P_u \| M),
\quad M = \frac{P_v + P_u}{2}
\end{equation}

\textbf{Giả thuyết:} Các TAZ có nhãn vùng rời rạc (ví dụ: công nghiệp và trường học)
sẽ có JSD cao hơn so với các TAZ có vùng chồng lấp.

\textbf{Kết quả thực chứng:} Kết quả phân tích thống kê trên tập dữ liệu HCM-Zone xác nhận giả thuyết này.
Phân kỳ JSD trung bình giữa các cặp TAZ có phân bổ chức năng tương đồng (Cosine Similarity $> 0.5$) chỉ là $0.0108$,
trong khi phân kỳ JSD trung bình giữa các cặp TAZ khác biệt chức năng hoàn toàn (Cosine Similarity $= 0$) lên tới $0.0177$ (tăng khoảng $64\%$).
Điều này chứng minh tính không đồng nhất về phân phối dữ liệu (Non-IID) giữa các vùng chức năng đô thị khác nhau.

%% ============================================================
\section{Kế Hoạch Triển Khai}
%% ============================================================

\begin{enumerate}
    \item \textbf{Tuần 1}: Chạy các truy vấn Overpass API cho tất cả 17 TAZ → \texttt{node\_zone\_labels.csv}
    \item \textbf{Tuần 1--2}: Thu thập dữ liệu TomTom bằng \texttt{tomtom\_collector.py}
          trong ít nhất 7 ngày (ghi lại toàn bộ mẫu hình hàng tuần bao gồm cả giờ cao điểm)
    \item \textbf{Tuần 2}: Xây dựng ma trận kề từ OSRM, chuyển đổi các cạnh TomTom thành
          đặc trưng TAZ, tính toán các số liệu thống kê động lực JSD
    \item \textbf{Tuần 3}: Triển khai và huấn luyện Zone-Aware AH-GNN, chạy các nghiên cứu cắt bỏ
    \item \textbf{Tuần 4}: Viết bài báo (Từ phần Giới thiệu đến Thực nghiệm)
\end{enumerate}

%% ============================================================
\section{Đóng Góp Dự Kiến}
%% ============================================================

\begin{enumerate}
    \item \textbf{Định dạng bài toán}: định nghĩa chính thức đầu tiên về đồ thị giao thông
          nhận biết vùng đa nhãn với chứng minh sự thiếu hụt của ma trận kề.
    \item \textbf{Kiến trúc}: Nhúng Vùng + Sinh Trọng Số Điều Biến Theo Vùng
          + Ma Trận Kề Thích Ứng Thiên Lệch Vùng.
    \item \textbf{Tập dữ liệu}: HCM-Zone — tập dữ liệu công khai đầu tiên kết hợp định tuyến OSRM,
          giao thông thời gian thực TomTom, và các chú thích vùng đa nhãn OSM cho
          TP. Hồ Chí Minh.
    \item \textbf{Thực nghiệm}: đánh giá phân lớp theo vùng thể hiện những
          cải thiện đặc biệt tại các TAZ đa vùng.
\end{enumerate}

%% ============================================================
\begin{thebibliography}{99}
\bibitem{li2018diffusion} Li, Y., et al. (2018). Diffusion Convolutional Recurrent Neural
Network. \textit{ICLR}.
\bibitem{wu2019graph} Wu, Z., et al. (2019). Graph WaveNet. \textit{IJCAI}.
\bibitem{yu2018spatio} Yu, B., et al. (2018). Spatio-Temporal Graph Convolutional Networks.
\textit{IJCAI}.
\bibitem{guo2021learning} Guo, S., et al. (2021). Learning Dynamics and Heterogeneity of
Spatial-Temporal Graph Data for Traffic Forecasting. \textit{IEEE TKDE}.
\bibitem{wang2019heterogeneous} Wang, X., et al. (2019). Heterogeneous Graph Attention
Network. \textit{WWW}.
\bibitem{wang2021self} Wang, X., et al. (2021). Self-supervised Heterogeneous Graph
Pre-training. \textit{NeurIPS}.
\bibitem{lv2021we} Lv, Q., et al. (2021). Are We Really Making Much Progress? \textit{KDD}.
\bibitem{wang2020urban2vec} Wang, Z., et al. (2020). Urban2Vec. \textit{AAAI}.
\end{thebibliography}

\end{document}
